\chapter{Toenail infection}
\index{Toenail infection}

\section*{Definition and Overview}
Toenail infections represent a group of pathological conditions affecting the nail unit of the feet. This anatomical unit includes the nail plate, the underlying nail bed, the nail matrix where nail growth begins, and the surrounding periungual skin folds. The most common toenail infection is onychomycosis, also known as tinea unguium, which is fungal in origin. Additionally, bacterial infections such as acute and chronic paronychia frequently affect the surrounding nail folds, sometimes extending subungual.

While historically considered a cosmetic issue, toenail infections are now recognised as significant medical conditions. They cause physical discomfort, pain, functional limitations, and significant psychosocial distress. The toenail acts as a protective shield for the distal digit and plays a key role in gait and balance. Destruction of the nail plate due to infection can interfere with walking and stability.

In patients with systemic disorders like diabetes mellitus, peripheral arterial disease, or chronic venous insufficiency, or in those who are immunocompromised, toenail infections can lead to serious secondary bacterial infections, cellulitis, osteomyelitis, and gangrene. Prompt and accurate diagnosis, followed by targeted treatment, is necessary to prevent these complications.

\section*{Epidemiology}
Onychomycosis is the most common nail disease, accounting for approximately 50\% of all nail disorders. The global prevalence in the general population is estimated between 10\% and 12\%. The epidemiology of toenail infections is strongly influenced by age, gender, lifestyle, and underlying health conditions.

The prevalence of toenail infections increases with age. It is rare in children (affecting less than 1\%) but rises to approximately 20\% in individuals over 60, and exceeds 50\% in those over 70. This increase is driven by slower nail growth, reduced peripheral blood flow, cumulative trauma, and a higher prevalence of systemic comorbidities.

Epidemiological studies show that adult males are affected by dermatophytic onychomycosis at a rate 1.5 to 3 times higher than females, likely due to occupational exposures, athletic participation, and regular use of occlusive footwear. Conversely, yeast infections and chronic paronychia are more common in females, often linked to frequent foot immersion in water.

Specific high-risk groups include patients with diabetes mellitus, where prevalence reaches 26\% to 30\% due to neuropathy and microvascular disease. Athletes, particularly runners and swimmers, show prevalence rates exceeding 25\% due to microtrauma and exposure to warm, humid communal spaces. Patients with psoriasis also have higher rates of secondary nail infections.

\section*{Aetiology and Pathophysiology}
Toenail infections arise from interactions between pathogenic organisms, host susceptibility, and environmental factors. Fungal pathogens cause most cases of onychomycosis, while bacteria typically cause acute periungual infections.

The primary causative pathogens in toenail infections include:
\begin{itemize}
    \item \textbf{Dermatophytes:} These filamentous, keratinophilic fungi cause 80\% to 90\% of all onychomycosis cases. \emph{Trichophyton rubrum} is the most common, followed by \emph{Trichophyton mentagrophytes} (specifically var. \emph{interdigitale}). They produce keratinases, enzymes that degrade nail keratin, allowing the fungi to invade the nail unit.
    \item \textbf{Non-Dermatophyte Moulds (NDMs):} Accounting for 10\% to 15\% of infections, these include \emph{Neoscytalidium dimidiatum}, \emph{Fusarium} species, \emph{Aspergillus} species, and \emph{Scopulariopsis brevicaulis}. They are opportunistic pathogens that infect nail plates already damaged by trauma or disease.
    \item \textbf{Yeasts:} Fungi of the genus \emph{Candida}, particularly \emph{Candida albicans} and \emph{Candida parapsilosis}, cause 5\% to 10\% of toenail infections. They are typically associated with fingernail infections but can affect toenails in immunocompromised patients or those exposed to chronic moisture.
    \item \textbf{Bacterial Pathogens:} Acute paronychia of the toenail is caused by pyogenic bacteria, mainly \emph{Staphylococcus aureus} or \emph{Streptococcus pyogenes}, entering through cuts, hangnails, or ingrown toenails. Additionally, \emph{Pseudomonas aeruginosa} can colonise the subungual space of a detached nail, causing "green nail syndrome" (chloronychia).
\end{itemize}

The pathophysiology of onychomycosis involves:
\begin{itemize}
    \item \textbf{Invasion:} Fungal spores germinate and hyphae penetrate the stratum corneum of the nail bed and the ventral surface of the nail plate.
    \item \textbf{Subungual Hyperkeratosis:} In response to fungal invasion, the nail bed produces excess keratin debris. This lifts the nail plate (onycholysis) and causes it to thicken.
    \item \textbf{Destruction:} As the fungus spreads toward the nail matrix, the nail plate becomes structurally defective, turning brittle, discoloured, and prone to crumbling.
    \item \textbf{Immunological Evasion:} Because the nail plate is avascular, the fungi are shielded from circulating immune cells and systemic antibodies, leading to a chronic infection.
\end{itemize}

Predisposing risk factors include:
\begin{itemize}
    \item \textbf{Microtrauma:} Repetitive mechanical pressure from footwear or sports damages the nail attachment and creates entry points for fungi.
    \item \textbf{Occlusive Footwear:} Tight, non-breathable shoes create a warm, moist microenvironment that promotes fungal growth.
    \item \textbf{Impaired Circulation:} Peripheral arterial disease (PAD) and chronic venous insufficiency reduce the delivery of oxygen and immune cells to the toes.
    \item \textbf{Neuropathy:} Diabetic sensory neuropathy reduces the patient's awareness of minor nail trauma, allowing infections to progress.
    \item \textbf{Genetic Predisposition:} An autosomal dominant pattern of inheritance may increase susceptibility to \emph{Trichophyton rubrum} infections.
\end{itemize}

\section*{Clinical Features}
The clinical presentation depends on the pathogen, the duration of the infection, and the host's anatomical response. Fungal toenail infections are classified into five major clinical subtypes.

The clinical subtypes and their features include:
\begin{itemize}
    \item \textbf{Distal Lateral Subungual Onychomycosis (DLSO):} The most common form, starting at the distal edge or lateral borders of the nail and spreading proximally. Signs include subungual hyperkeratosis (yellow-white keratin debris under the nail), onycholysis (separation of the nail from the bed), and yellow-brown streaks.
    \item \textbf{Superficial White Onychomycosis (SWO):} Characterised by chalky-white patches or striations on the dorsal surface of the nail plate. Fungal hyphae invade the superficial nail without involving the nail bed. The affected area is soft and powdery, easily scraped off.
    \item \textbf{Proximal Subungual Onychomycosis (PSO):} The fungus invades through the proximal nail fold and migrates distally, appearing as a white-yellow discolouration near the lunula. This subtype is a clinical marker of immunodeficiency, particularly HIV.
    \item \textbf{Endonyx Onychomycosis:} The fungus directly invades the interior of the nail plate. It presents with a milky-white discolouration but lacks subungual hyperkeratosis and onycholysis.
    \item \textbf{Total Dystrophic Onychomycosis (TDO):} The end-stage of chronic infection where the entire nail plate is destroyed, leaving a thick, yellow-brown, deformed, and crumbly mass of keratin debris.
\end{itemize}

Bacterial toenail infections present distinct clinical features:
\begin{itemize}
    \item \textbf{Acute Paronychia:} Sudden onset of pain, redness, swelling, and warmth along the lateral or proximal nail folds, often with a localized pus abscess.
    \item \textbf{Green Nail Syndrome (Chloronychia):} A greenish-black discolouration of the nail plate caused by \emph{Pseudomonas aeruginosa} colonising the subungual space under a detached nail plate.
\end{itemize}

Patients frequently report secondary symptoms, including pain when wearing shoes or walking, difficulty trimming nails, and significant embarrassment about the cosmetic appearance of their feet.

\section*{Differential Diagnosis}
Because approximately 50\% of all nail abnormalities are non-infectious, visual diagnosis alone is insufficient. A precise differential diagnosis prevents unnecessary treatment.

Key differential diagnoses include:
\begin{itemize}
    \item \textbf{Nail Psoriasis:} Mimics onychomycosis with features like nail pitting (small depressions), "oil drop" discolouration under the nail, subungual hyperkeratosis, and splinter haemorrhages. It is often accompanied by skin psoriasis and affects multiple nails symmetrically.
    \item \textbf{Lichen Planus:} An inflammatory condition causing longitudinal ridging, thinning of the nail plate, splitting, and pterygium formation (fusion of the nail fold to the nail bed), which can lead to permanent nail loss.
    \item \textbf{Traumatic Dystrophy:} Caused by acute injury or chronic microtrauma. It presents with subungual haematomas, horizontal Beau's lines, onycholysis, and thickening. It is typically confined to the injured toe and lacks progressive fungal debris.
    \item \textbf{Subungual Melanoma:} A rare, life-threatening malignancy presenting as a dark brown or black longitudinal band (melanonychia) that expands. The extension of pigment to the surrounding skin (Hutchinson's sign) warrants an urgent biopsy.
    \item \textbf{Onychogryphosis:} Characterised by extreme thickening and distortion of the nail plate, which curves in a horn-like shape. Commonly seen in elderly patients due to neglect or poor circulation, without active fungal infection.
    \item \textbf{Yellow Nail Syndrome:} A rare systemic disorder featuring slow-growing, thickened, pale yellow nails without cuticles or lunulae, associated with lymphoedema and respiratory disease.
    \item \textbf{Contact Dermatitis and Eczema:} Inflammatory skin diseases causing irregular transverse ridging, scaling, and erythema of the surrounding skin.
\end{itemize}

\section*{When to Seek Medical Care}
While many toenail infections develop slowly, certain signs demand prompt medical attention to avoid severe complications.

Consult a healthcare professional (such as a GP, podiatrist, or dermatologist) if you experience:
\begin{itemize}
    \item Any signs of toenail infection if you have diabetes mellitus, peripheral neuropathy, or peripheral arterial disease. In these groups, minor nail trauma can lead to non-healing ulcers, osteomyelitis, or tissue necrosis.
    \item Spreading redness, swelling, warmth, or severe pain in the skin surrounding the toe, which indicates cellulitis.
    \item Red streaks extending up the foot or leg, indicating a spreading infection (lymphangitis).
    \item Pus draining from the nail folds, indicating a localized abscess.
    \item A dark brown or black spot or band under the nail that is expanding or not caused by a recent injury.
    \item Signs of infection if you are immunocompromised due to chemotherapy or disease.
    \item Pain that interferes with walking, standing, or wearing shoes.
\end{itemize}

Seek emergency medical services immediately if you experience systemic symptoms of severe infection, such as high fever, chills, rapid heart rate, or confusion. Contact emergency services:
\begin{itemize}
    \item \textbf{local emergency services} in many countries.
    \item \textbf{911} in the United States and Canada.
    \item \textbf{999} or \textbf{112} in the United Kingdom and Europe.
\end{itemize}

\section*{Diagnosis and Investigations}
Visual inspection has an accuracy of only 50\% to 65\%. Laboratory confirmation is required before starting oral antifungal therapy due to potential side effects and drug interactions.

Diagnostic investigations include:
\begin{itemize}
    \item \textbf{Clinical Examination:} Assessing the pattern of nail changes, the presence of subungual debris, discolouration, and checking for skin infections (such as tinea pedis).
    \item \textbf{Potassium Hydroxide (KOH) Preparation:} A rapid, in-office test. Subungual debris and nail clippings are treated with 10\% to 20\% KOH to dissolve keratin, allowing the clinician to visualise fungal hyphae or budding yeast under a microscope. Sensitivity is 60\% to 80\%; it does not identify the species.
    \item \textbf{Fungal Culture:} Samples are grown on selective culture media (e.g., Sabouraud dextrose agar) to identify the specific pathogen. Culturing helps differentiate dermatophytes, yeasts, and moulds, but it takes 2 to 6 weeks and has a false-negative rate of 30\% to 50\%.
    \item \textbf{Histopathology with Periodic Acid-Schiff (PAS) Staining:} The most sensitive test (sensitivity >90\%) for confirming fungal presence. A nail clipping is processed and stained with PAS, which highlights fungal cell walls in bright magenta. It does not identify the specific species.
    \item \textbf{Polymerase Chain Reaction (PCR) Assays:} Molecular testing that detects fungal DNA directly from nail clippings within 24 to 48 hours. It is highly sensitive and specific, identifying the exact pathogen, but is more expensive.
\end{itemize}

\section*{Management}
Treatment requires patience and adherence, as toenails grow slowly, taking 12 to 18 months to fully replace themselves. The choice of therapy depends on the severity, causative organism, drug interactions, and patient health.

Management options include:
\begin{itemize}
    \item \textbf{Oral Antifungals:} The most effective treatments for moderate-to-severe onychomycosis. Terbinafine (250 mg daily for 12 weeks) is first-line for dermatophytes, achieving a 70\% to 80\% mycological cure. Itraconazole (200 mg daily for 12 weeks, or pulse therapy) is preferred for yeast and moulds. Fluconazole (150 mg to 300 mg once weekly) is an alternative. Because oral antifungals can cause liver toxicity, baseline and periodic liver function tests (LFTs) are recommended.
    \item \textbf{Topical Antifungals:} Suitable for mild-to-moderate infections without matrix involvement (less than 50\% nail involvement) or when oral drugs are contraindicated. Options include Amorolfine 5\% nail lacquer (once or twice weekly), Ciclopirox 8\% lacquer (daily), Efinaconazole 10\% solution (daily), and Tavaborole 5\% solution (daily). Cure rates are lower (15\% to 30\%) because of poor penetration through the nail plate, but there are no systemic side effects.
    \item \textbf{Combination Therapy:} Combining oral terbinafine with a topical lacquer (like amorolfine) improves cure rates in severe or resistant cases compared to using oral therapy alone.
    \item \textbf{Laser and Physical Therapies:} Nd:YAG and diode lasers are approved for temporarily clearing nail discolouration. They use photothermal energy to inhibit fungal growth. Clinical studies show variable long-term cure rates, and lasers are typically used as adjunctive options.
    \item \textbf{Debridement and Surgical Avulsion:} Chemical debridement with a 40\% urea ointment under occlusion dissolves thick nail tissue, helping topicals penetrate. In severe, painful cases, surgical removal (avulsion) of the nail plate under local anaesthesia may be performed.
    \item \textbf{Bacterial Infection Management:} Acute paronychia is managed with warm compresses, elevation, and, if necessary, incision and drainage. Oral antibiotics (such as cephalexin) are prescribed for cellulitis. Green nail syndrome is managed by keeping the nail dry, clipping the detached portion, and applying topical dilute acetic acid or antibiotic drops.
\end{itemize}

\section*{Complications}
Untreated toenail infections can lead to significant local and systemic complications, especially in vulnerable patient groups.

Complications include:
\begin{itemize}
    \item \textbf{Cellulitis:} Fungal or bacterial nail damage can compromise the skin barrier, allowing bacteria to enter and cause spreading soft-tissue infections in the foot or lower leg.
    \item \textbf{Diabetic Foot Complications:} Diabetic neuropathy can prevent patients from feeling pain from thickened, deformed nails. This mechanical pressure can lead to subungual ulcers, secondary bacterial infections, osteomyelitis, and gangrene, which may require amputation.
    \item \textbf{Permanent Nail Dystrophy:} Long-standing infection can damage the nail matrix, causing chronic abnormal nail growth, thickening, and permanent cosmetic deformity.
    \item \textbf{Chronic Pain and Mobility Impairment:} Painful, thick nails make walking and standing uncomfortable. This can alter gait patterns, leading to secondary joint pain and increasing the risk of falls in the elderly.
    \item \textbf{Antifungal Toxicity:} Oral therapies carry risk. Terbinafine and itraconazole can cause drug-induced liver injury, ranging from elevated liver enzymes to liver failure, and interact with common drugs like beta-blockers and statins.
    \item \textbf{Psychosocial Distress:} Cosmetic deformity of the nails can cause significant embarrassment, leading to anxiety and social withdrawal.
\end{itemize}

\section*{Prognosis and Outcomes}
The prognosis depends on the patient's age, the causative organism, the severity of the infection, and underlying health conditions. Outcomes are measured by mycological cure (negative tests) and clinical cure (a completely normal-looking nail).

Oral terbinafine is the most successful therapy, achieving mycological cure rates of 70\% to 80\% and clinical cure rates of 35\% to 50\%. A nail may be mycologically cured but remain permanently damaged due to scarring of the nail bed or matrix. Topical treatments alone have lower mycological cure rates of 15\% to 35\%.

Recurrence and reinfection are common, occurring in 20\% to 50\% of patients within 3 years of treatment. This is due to ongoing fungal exposure, genetic susceptibility, incomplete clearance of the fungus, and poor adherence to prevention.

Factors that increase the risk of treatment failure include:
\begin{itemize}
    \item Severe nail dystrophy involving more than 50\% of the nail.
    \item Involvement of the nail matrix.
    \item The presence of a "dermatophytoma" (a dense mass of fungal hyphae in the nail).
    \item Poor compliance with the long treatment course.
    \item Co-existing diabetes or peripheral arterial disease.
    \item Advanced age (slower nail growth and poorer circulation).
\end{itemize}

\section*{Prevention and Self-Care}
Preventing toenail infections and avoiding recurrence requires consistent foot hygiene and simple lifestyle adjustments to minimise fungal exposure.

Preventive measures include:
\begin{itemize}
    \item \textbf{Foot Hygiene:} Wash feet daily with soap and water, and dry them thoroughly, especially between the toes, before putting on socks and shoes.
    \item \textbf{Socks and Footwear:} Wear moisture-wicking socks made of synthetic fibres or merino wool. Avoid 100\% cotton socks. Change socks immediately if they become damp. Choose breathable shoes (made of leather or canvas) and avoid tight shoes.
    \item \textbf{Shoe Care:} Rotate shoes daily to allow them to dry for 24 hours. Apply antifungal powders or sprays inside shoes daily.
    \item \textbf{Public Areas:} Wear protective footwear (sandals or flip-flops) in communal wet areas like public showers, locker rooms, and pool decks.
    \item \textbf{Nail Care:} Cut toenails straight across and file down thick areas. Never use the same clippers on infected and healthy nails, and sterilise nail tools with rubbing alcohol after use.
    \item \textbf{Do Not Share:} Avoid sharing socks, shoes, towels, or nail tools.
    \item \textbf{Early Treatment:} Treat athlete's foot (tinea pedis) immediately with over-the-counter creams to prevent it from spreading to the nails.
    \item \textbf{Nail Salons:} Ensure any nail salon you visit sterilises instruments in an autoclave, or bring your own tools.
\end{itemize}

\section*{Sources and Further Reading}
For additional information and clinical guidelines, consult these professional organizations:
\begin{itemize}
    \item \textbf{American Academy of Dermatology (AAD):} Patient education and clinical guides on managing fungal nail infections. URL: https://www.aad.org/public/diseases/a-z/nail-fungus
    \item \textbf{Centers for Disease Control and Prevention (CDC):} Detailed overview of fungal nail infections, clinical guidelines, and prevention. URL: https://www.cdc.gov/fungal/diseases/fungal-nail-infections
    \item \textbf{British Association of Dermatologists (BAD):} Clinical guidelines and patient information sheets on onychomycosis. URL: https://www.bad.org.uk/patient-information-leaflets
    \item \textbf{DermNet New Zealand:} Peer-reviewed dermatological database with clinical images and treatment protocols for nail infections. URL: https://dermnetnz.org/topics/onychomycosis
    \item \textbf{Mayo Clinic:} Clinical diagnostic and treatment guides for nail fungus. URL: https://www.mayoclinic.org/diseases-conditions/nail-fungus
    \item \textbf{Harvard Health Publishing:} Educational articles discussing foot hygiene, prevention, and risks of untreated nail infections. URL: https://www.health.harvard.edu/a\_to\_z/fungal-nail-infections
\end{itemize}